\begin{table}[htbp]
\centering
\caption{Pre-trends test. Four leads of the exposure interaction are added to the main specification. Under parallel trends all four are zero. Standard errors, clustered by county, in parentheses. $^{***}p<0.01$, $^{**}p<0.05$, $^{*}p<0.10$.}
\label{tab:placebo}
\begin{tabular}{lcc}
\toprule
 & Full panel & Excluding \\
 & 2002--2019 & 2009--2012 \\
 & (1) & (2) \\
\midrule
Exposure $\times\ \Delta$ federal funds rate & 0.0252$^{***}$ & 0.0229$^{***}$ \\
 & (0.0060) & (0.0057) \\
\quad lead 1 & 0.0157$^{***}$ & 0.0146$^{***}$ \\
 & (0.0052) & (0.0052) \\
\quad lead 2 & 0.0172$^{***}$ & 0.0157$^{***}$ \\
 & (0.0055) & (0.0055) \\
\quad lead 3 & -0.0217$^{***}$ & -0.0226$^{***}$ \\
 & (0.0043) & (0.0044) \\
\quad lead 4 & -0.0030 & -0.0019 \\
 & (0.0097) & (0.0098) \\
\midrule
County fixed effects & Yes & Yes \\
Quarter fixed effects & Yes & Yes \\
Observations & 212,468 & 162,484 \\
\bottomrule
\end{tabular}

\begin{minipage}{0.92\textwidth}\vspace{4pt}\footnotesize Leads one, two and three reject zero at five percent in both windows, so the parallel-trends assumption fails and the difference-in-differences estimate cannot be read causally. The paper reports a count of two, the leads that also reject under randomization inference; see Section~\ref{sec:pretrends}.\end{minipage}
\end{table}

